\documentclass[12pt]{article}

\usepackage[utf8]{inputenc}
\usepackage[spanish]{babel}
\usepackage[shortlabels]{enumitem}

\usepackage{mathtools}
\usepackage{amssymb}
\usepackage{amsthm}

\usepackage{geometry}
\geometry{margin=2.5cm}

\usepackage{hyperref}

\title{Ejercicios Limites Y Continuidad}
\author{}

\begin{document}
\maketitle
\section*{Ejercicio 1}
Calcular el limite
\[
\lim_{x\to 0}\frac{\sqrt{\,1+\ln(1+x)\,}-\sqrt{\,1-\ln(1+x^2)\,}}{x}.
\]
\section*{Ejercicio 2}
Calcular el limite
\[
\lim_{x\to 0}\frac{1+2(\cos x)^3-3\sqrt{\cos(2x)}}{(\sin x)^2}.
\]
\section*{Ejercicio 3}
Si $f(x):(-\infty,2]\to \mathbb{R}$ es la funcion definida por
\[
f(x)=
\begin{cases}
\dfrac{a}{1+e^{-bx}}, & x\in(-\infty,0),\\[8pt]
\dfrac{3x+2}{x^2-x-6}, & x\in[0,2],
\end{cases}
\qquad a,b\in\mathbb{R},
\]
determinar de forma razonada para que valores de $a$ y $b$ la funcion es continua en todo su dominio.
\section*{Ejercicio 4}
Sea la funcion real de variable real definida a continuacion:
\[
f(x)=
\begin{cases}
\dfrac{\sen(x)}{1+e^{1/\sen(x)}}, & x\ne 0,\\[8pt]
0, & x=0.
\end{cases}
\]
\begin{enumerate}[a)]
\item Estudiar la continuidad de $f(x)$ en $x=0$, $x=\pi/2$ y $x=\pi$. En caso de discontinuidad, indicar claramente el tipo.
\end{enumerate}
\section*{Ejercicio 5}
Calcular el limite
\[
\lim_{x\to 0}\frac{\sqrt{1+x\sen(x)}-\sqrt{\cos(2x)}}{\tan^2\!\left(\frac{x}{2}\right)}.
\]
\section*{Ejercicio 6}
Calcular el siguiente limite:
\[
\lim_{x\to \pi/2}\left(\frac{e^x\cos(x)+1}{e^{\cos(x)}}\right)^{1/\cos(x)}.
\]
\section*{Ejercicio 7}
Dada la siguiente funcion $f:(-\pi,\pi)\to \mathbb{R}$, determinar de forma razonada el valor $\alpha\in\mathbb{R}$ para que la funcion sea continua en $x=0$. En caso de no existir dicho valor, justificar la respuesta.
\[
f(x)=
\begin{cases}
\dfrac{\ln\!\left(\dfrac{\sen(x)}{x}\right)}{x^2}, & x\in(-\pi,0)\cup(0,\pi),\\[10pt]
\alpha, & x=0.
\end{cases}
\]
\section*{Ejercicio 8}
Dada la funcion
\[
f(x)=
\begin{cases}
\dfrac{\sqrt{1+kx}-\sqrt{1-kx}}{x}, & -1\le x<0,\\[10pt]
\dfrac{2x+1}{x-1}, & 0\le x\le 1,
\end{cases}
\]
determinar los valores del parametro $k\in\mathbb{R}$ para los cuales la funcion es continua en $x=0$.
\end{document}